% !PRINTED 462
\weberpage{462}
%<<CONT
$\alpha, \alpha_1, \alpha_2 \ldots \alpha_{n-1}$ of $f(x) = 0$ vanish, i.\,e. we have
\[
\varphi(\Theta) = 0,\quad \varphi(\Theta_1) = 0,\quad \varphi(\Theta_2) = 0 \ldots \varphi(\Theta_{n-1}) = 0.
\]

If, then, the values conjugate to $\Theta$ are all different from
one another, $\varphi(t)$ is identical with $\Phi(t)$, i.\,e.
$\Phi(t)$ is irreducible. But if among the values conjugate to
$\Theta$ only $n_1$ different ones are present,
say $\Theta, \Theta_1, \Theta_2 \ldots \Theta_{n_1-1}$, then
\weq{10}{\varphi(t) = (t - \Theta)(t - \Theta_1) \ldots (t - \Theta_{n_1-1}),}
and every other irreducible factor of $\Phi(t)$, since it must
vanish for at least one of the conjugate values $\Theta$ and
consequently for all of them, is identical with $\varphi(t)$. Hence
$\Phi(t)$ is a power of $\varphi(t)$, say
\weq{11}{\Phi(t) = \varphi(t)^{n_2}}
and
\weq{12}{n = n_1 n_2.}
From this, then, follows the theorem:

\gesperrt{2. The function $\Phi(t)$ is either irreducible or it
is a power of an irreducible function. The $n$ numbers conjugate
to a number in $\Omega(\alpha)$ are either all different from
one another, or they fall into $n_1$ systems of
$n_2$ numbers each, equal among themselves. In the first case $\Phi(t)$
is irreducible, in the second it is the $n_2$th power of an
irreducible function of degree $n_1$ in $\Omega$.}

A quantity $\Theta$ in $\Omega(\alpha)$ that is different from all its
conjugate numbers, and that therefore satisfies an
irreducible equation of degree $n$, is called a
\gesperrt{primitive quantity of the field}. By the theorem \S.\,143, 1. infinitely
many such primitive quantities can be determined, indeed in such a way that
the coefficients of $\Theta = \chi(\alpha)$ are rational numbers.
One need only dispose of the coefficients of $\chi$ in such a way
that among the conjugate quantities $\chi(\alpha)$ no equal ones
occur.

\gesperrt{3. Every quantity $\omega$ of the field $\Omega(\alpha)$ can
be expressed rationally in terms of an arbitrary primitive quantity $\Theta$
of the field.}

For if $\omega, \omega_1, \omega_2 \ldots \omega_{n-1}$ are the numbers
conjugate to $\omega$, and likewise $\Theta, \Theta_1, \Theta_2
\ldots \Theta_{n-1}$ those conjugate to $\Theta$, and
%CONT>>
